\section{Traducción en términos del grupo fundamental}
\label{IX.5}

Sea
$$
g\colon S'\to S
$$
un \emph{morfismo de descenso efectivo} para la categoría fibrada de los
\emph{recubrimientos étales} de preesquemas; por ejemplo, un morfismo propio,
sobreyectivo y de presentación finita~\eqref{IX.4.12},
o un morfismo fielmente plano y cuasicompacto. Introduzcamos, como de
costumbre,~$S''$ y~$S'''$, y denotemos por~$\cal{C}$,~$\cal{C}'$,
~$\cal{C}''$ y~$\cal{C}'''$ las categorías de los recubrimientos étales de
~$S$,~$S'$,~$S''$ y~$S'''$, respectivamente. Se tiene entonces un diagrama
$2$-exacto de categorías
\begin{equation*}
\label{eq:IX.5.*}
\tag{$*$}
\xymatrix@C=.5cm{
{\cal{C}} \ar[r]^-{~p^*} & {\cal{C}'}
\ar@/^2pt/[rr]^{p_1^*,p_2^*} \ar@/_2pt/[rr] && {\cal{C}''}
\ar@/^4pt/[rrr]^{p_{21}^*,p_{32}^*,p_{31}^*} \ar[rrr] \ar@/_4pt/[rrr] &&&
{\cal{C}'''}
}
\end{equation*}
correspondiente al diagrama
$$
\xymatrix@C=.5cm{ S & \ar[l]_-{~p} S' && \ar@/^2pt/[ll] \ar@/_2pt/[ll]_-{p_1,p_2}
S'' &&& \ar@/^4pt/[lll] \ar[lll] \ar@/_4pt/[lll]_{p_{21},p_{32},p_{31}} S'''.
}
$$
Supongamos
\marginpar{243}
que los preesquemas~$S$,~$S'$,~$S''$ y~$S'''$ son sumas disjuntas de
preesquemas conexos, como ocurre, en particular, si son localmente conexos y,
con mayor razón, si son localmente noetherianos (por ejemplo, si~$S'$ es de
tipo finito sobre~$S$ localmente noetheriano). Las categorías~$\cal{C}$,
~$\cal{C}'$, …, de~\eqref{eq:IX.5.*} son entonces categorías
multigaloisianas~\Ref{V.9}, descritas cada una por una familia de grupos
topológicos compactos totalmente disconexos, a saber, los grupos
fundamentales de las componentes conexas de los preesquemas~$S$,~$S'$,
~$S''$ y~$S'''$. Para simplificar, supondremos que~$S$ es conexo y daremos un
procedimiento de cálculo de su grupo fundamental en términos de la categoría
fibrada formada por~$\cal{C}'$,~$\cal{C}''$ y~$\cal{C}'''$, explicitada
convenientemente mediante los grupos fundamentales que describen estas
categorías. El lector observará que el procedimiento esbozado vale, de hecho,
en el marco general de las categorías multigaloisianas (que no tienen por qué
proceder de preesquemas dados~$S$,~$S'$,~$S''$ y~$S'''$). Es, por lo demás,
el análogo del procedimiento bien conocido para calcular el grupo fundamental
de un espacio topológico~$S$, unión localmente finita de subespacios
cerrados~$S_i$ (o unión arbitraria de subespacios abiertos~$S_i$), mediante
los grupos fundamentales de las componentes
\ifthenelse{\boolean{orig}}
{de los~$X_i$}
{de los~$S_i$}
y de las componentes de los~$S_i\cap S_j$. Naturalmente, la situación análoga
para los preesquemas cae dentro del marco general del descenso, introduciendo
el preesquema~$S'$, suma de los~$S_i$, y el morfismo canónico
$g\colon S'\to S$.

Pongamos
$$
E'=\pi_0(S'), \quad E''=\pi_0(S''), \quad E'''=\pi_0(S'''),
$$
donde~$\pi_0$ designa el funtor «conjunto de las componentes conexas». Como
los productos fibrados de~$S'$ sobre~$S$ forman un objeto simplicial de
~$\Sch$, el funtor~$\pi_0$ lo transforma en un conjunto simplicial del cual
~$E'$,~$E''$ y~$E'''$ son las componentes de dimensiones~$0$,~$1$ y~$2$.
Habremos
de utilizar las aplicaciones simpliciales
$$
q_i=\pi_0(p_i),\quad(i=1,2)\quad\text{y}
\quad q_{ij}=\pi_0(p_{ij}),\quad (i,j)=(2,1),(3,2),(3,1),
$$
puestas de manifiesto en el diagrama
\begin{equation*}
\label{eq:IX.5.1}
\tag{1} {\xymatrix@C=.5cm{ E' && \ar@/^2pt/[ll] \ar@/_2pt/[ll]_{q_1,q_2} E'' &&&
\ar@/^4pt/[lll] \ar[lll] \ar@/_4pt/[lll]_{q_{21},q_{32},q_{31}} E'''. }}
\end{equation*}
Los elementos
\marginpar{244}
de~$E'$ se denotarán con una prima, como~$s'$; los de~$E''$ \resp los de
~$E'''$ se denotarán con~$''$ \resp con~$'''$. Que~$S$ sea conexo se expresa por
$\pi_0(K)=0$, donde~$K$ es el conjunto simplicial definido por
$g\colon S'\to S$, o también por el hecho de que la relación de equivalencia
en~$E'$ engendrada por el par de aplicaciones~$(q_1,q_2)$ es transitiva.

Elegiremos, de una vez por todas, un elemento~$s_0'$ de~$E'$ y, para cada
~$s'$ de~$E'$, un elemento~$\overline{s'}\in E''$ tal que\footnote{Debe
advertirse que el elemento~$\overline{s'}$, cuya existencia se admite
implícitamente, no existe en todos los casos. Por consiguiente, el
teorema~\Ref{IX.5.1}, tal como está enunciado, no se aplica en todos los casos.
Sin embargo, no es difícil, inspirándose en el texto, modificar su enunciado
de modo que proporcione un método de cálculo aplicable en todos los casos. En
particular, los corolarios de dicho teorema son válidos tal cual.}
$$
q_1(\overline{s'})=s_0', \quad q_2(\overline{s'})=s',
$$
poniendo así de manifiesto la conexidad de~$S$. Para cada~$s'\in E'$, elijamos
un punto geométrico~$\underline{s}'$ en la componente conexa~$s'$ de~$S'$;
este punto intervendrá, de hecho, mediante el funtor fibra~$F_{s'}'$ asociado
sobre la categoría multigaloisiana~$\cal{C}'$. El grupo de automorfismos de
este funtor, es decir, el grupo fundamental de~$S'$ en~$\underline{s}'$, se
denotará por~$\pi_{s'}$. Elijamos del mismo modo puntos~$\underline{s}''$ y
~$\underline{s}'''$, y por tanto funtores~$F_{s''}''$ y~$F_{s'''}'''$, de
donde se obtienen los grupos fundamentales~$\pi_{s''}$ y~$\pi_{s'''}$. Así,
$$
\pi_{s'}=\pi_1(S',\underline{s}'), \quad
\pi_{s''}=\pi_1(S'',\underline{s}''), \quad
\pi_{s'''}=\pi_1(S''',\underline{s}''').
$$
Para todo~$s''\in E''$, $p_1(\underline{s}'')$ pertenece a la misma
componente conexa que~$\underline{q_1(s'')}$; existe, pues, un isomorfismo de
funtores
\ifthenelse{\boolean{orig}}
{$F''_{s''}\circ p_1^* F'_{s'}$}
{$F''_{s''}\circ p_1^*\isomto F'_{q_1(s'')}$}
(es decir, una «clase de caminos» de~$p_1(\underline{s}'')$ a
~$\underline{q_1(s'')} $). La misma observación se aplica a~$q_2$ y a los
~$q_{ij}$. Elijamos todas estas clases de caminos:
$$
F''_{s''}\circ p_i^*\isomto F'_{q_i(s'')},\quad F_{s'''}'''\circ
p_{ij}^*\isomto
\ifthenelse{\boolean{orig}}{F'''_{q_{ij}(s''')}}{F''_{q_{ij}(s''')}},
$$
(para
\marginpar{245}
$i=1,2$ y $(i,j)=(2,1),(3,2),(3,1)$). Se obtienen, en particular,
homomorfismos de grupos
\begin{equation*}
\label{eq:IX.5.2}
\tag{2} {q_i^{s''}\colon \pi_{s''}\to\pi_{q_i(s'')},\quad
q_{ij}^{s'''}\colon \pi_{s'''}\to\pi_{q_{ij}(s''')},}
\end{equation*}
(con los mismos valores de~$i$ y~$(i,j)$). Recordemos, finalmente, que en la
estructura con clivaje de la categoría fibrada con fibras~$\cal{C}'$,
~$\cal{C''}$ y~$\cal{C'''}$ figuran también isomorfismos de funtores
$$
p_{21}^*p_1^*\isomto p_{31}^*p_1^*,\quad p_{21}^*p_2^*\isomto
p_{32}^*p_1^*,\quad p_{31}^*p_2^*\isomto p_{32}^*p_2^*,
$$
deducidos de sendos isomorfismos de ambos miembros con los~$u_i^*$
~($i=1,2,3$), donde los~$u_i$ son las tres proyecciones de~$S'''$ en~$S'$.
Al explicitar estos datos, para cada~$s'''$ se obtiene un elemento bien
determinado
\begin{equation*}
\label{eq:IX.5.3}
\tag{3} {a_i^{s'''}\in\pi_{v_i(s''')},}
\end{equation*}
(donde los~$v_i$, $i=1,2,3$, son las tres aplicaciones~$E'''\to E'$ definidas
por~$v_i=\pi_0(u_i)$), sometido además a las condiciones
$$
q_1^{s_1''}q_{21}^{s'''}=\mathrm{int}(a_1^{s'''})
q_1^{s_2''}q_{31}^{s'''}\qquad (s_1''=q_{21}(s'''),\
s_2''=q_{31}(s''')),
$$
y a las dos condiciones análogas, en las que intervienen~$a_2$ y~$a_3$. El
lector observará, por lo demás, que los datos~\eqref{eq:IX.5.1},
~\eqref{eq:IX.5.2} y~\eqref{eq:IX.5.3} permiten reconstruir, salvo equivalencia
de categorías fibradas, la categoría fibrada considerada, con fibras
~$\cal{C}'$,~$\cal{C}''$ y~$\cal{C}'''$. Deben permitir, pues, reconstruir en
principio~$\cal{C}$ salvo equivalencia y, por tanto, su grupo fundamental salvo
isomorfismo. Determinaremos, en realidad, el grupo fundamental en el punto
geométrico~$p(\underline{s}_0')$ de~$S$, es decir, el grupo de automorfismos
de~$F'_{s_0'}\circ p^*$.

Obsérvese que dar un objeto~$X'$ de~$\cal{C}'$ equivale esencialmente a dar
conjuntos finitos~$X'_{s'}$ ($s'\in E'$) sobre los que los~$\pi_{s'}$ actúan
continuamente. Un
\marginpar{246}
dato de pegado sobre tal objeto equivale entonces a dar, para cada
~$s''\in E''$, una biyección
$$
\varphi_{s''}\colon X'_{q_1(s'')}\isomto X'_{q_2(s'')}
$$
compatible con las acciones de~$\pi_{s''}$ sobre ambos miembros por medio de
los homomorfismos~$q_i^{s''}\colon \pi_{s''}\to\pi_{q_i(s'')}$. Tomando
primero los~$s''$ de la forma~$\overline{s'}$, se ve que tal dato define
biyeciones
$$
\psi_{s'}\colon X'_{s_0'}=F_0'(X')\to X_{s'}'
$$
lo que permite identificar todos los~$X'_{s'}$ con el mismo conjunto
~$F_0'(X')=X'_{s_0'}$, sobre el cual actuarán desde ese momento todos los
grupos~$\pi_{s'}$. Hecho esto, las biyecciones~$\varphi_{s''}$ corresponderán
a biyecciones
$$
g_{s''}\colon F_0'(X')\isomto F_0'(X'),
$$
sometidas, por una parte, a las relaciones de conmutación con~$\pi_{s''}$
\begin{equation*}
\label{eq:IX.5.a}
\tag*{a)} {g_{s''}q_1^{s''}(g'')=q_2^{s''}(g'')g_{s''}\qquad (s''\in E'',\ g''\in\pi_{s''}),}
\end{equation*}
y, por otra, a las relaciones
\begin{equation*}
\label{eq:IX.5.b}
\tag*{b)} {g_{\overline{s'}}=g_{\overline{s_0'}}\qquad (s'\in E'),}
\end{equation*}
que expresan cómo habíamos identificado entre sí los
~$\ifthenelse{\boolean{orig}}{X_{s'}}{X_{s'}'}$. Al explicitar
la condición para que tal dato de pegado sea, en efecto, un dato de descenso,
se obtienen las relaciones
\begin{equation*}
\label{eq:IX.5.c}
\tag*{c)} {a_3^{s'''}g_{q_{31}(s''')}a_1^{s'''}=
g_{q_{32}(s''')}a_2^{s'''}g_{q_{21}(s''')}\qquad (s'''\in E''').}
\end{equation*}

Esto proporciona una equivalencia entre la categoría de los objetos de
~$\cal{C}'$ provistos de un
\marginpar{247}
dato de descenso y la categoría de los conjuntos finitos sobre los que actúan
continuamente los grupos~$\pi_{s'}$, provistos además de biyecciones
~$g_{s''}$ que satisfacen las relaciones~\Ref{eq:IX.5.a},
~\Ref{eq:IX.5.b} y~\Ref{eq:IX.5.c}. Sea entonces~$G$ el grupo engendrado por
los grupos~$\pi_{s'}$ y los nuevos generadores~$g_{s''}$, sometidos a las
relaciones~\Ref{eq:IX.5.a},~\Ref{eq:IX.5.b} y~\Ref{eq:IX.5.c}, y sea~$\pi$
el grupo límite proyectivo de los cocientes de~$G$ por los subgrupos de índice
finito cuyas imágenes inversas en los grupos~$\pi_{s'}$ sean subgrupos
abiertos. También se dice que~$\pi$ es el \emph{grupo de tipo galoisiano
engendrado por los~$\pi_{s'}$ y los~$g_{s''}$, sometido a las
relaciones~\Ref{eq:IX.5.a},~\Ref{eq:IX.5.b} y~\Ref{eq:IX.5.c}}. Se comprueba
de inmediato que la categoría considerada es asimismo equivalente a la
categoría de conjuntos finitos sobre los que el grupo topológico~$\pi$ actúa
continuamente. Esto demuestra el enunciado siguiente:
